\subsection{A rate bracket must travel with the loan}
\label{sec:credit-rate-interval}

The preceding dated examples have a convenient property: their initial cash
flows identify an exactly representable effective rate. Ordinary schedules
need not do so. Consider an asset acquired on 1 January 2025 for 98 with
integral acquisition costs of 2. Expected receipts are 50 on 1 January 2026
and 60 on 1 January 2027. Its initial carrying amount is 100, including those
costs. For a supplied POCI classification, the original credit-adjusted rate
discounts these initial expected receipts; subsequent interest applies that
rate to amortised cost \citep[IFRS~9.5.4.1(a),B5.4.1--2,B5.4.7--8]{eu20231803}.

Retain the supplied ACT/365F convention and annual compounding. Both intervals
here are exactly 365 days. Write \(x=1+r\), where \(r>-1\) is the original
annual rate. The initial measurement requires
\[
 100=\frac{50}{x}+\frac{60}{x^2},\qquad
 100x^2-50x-60=0,\qquad
 x=\frac{5+\sqrt{265}}{20}.
\]
The positive solution gives a rate of about 6.3941\%. That decimal is useful
for reading the result; substituting it for the original rate would change
every later valuation. An implementation must carry a justified numerical
bound instead of presenting the rounded rate as exact.

More generally, let \(c_i\geq0\) be initial expected receipts at strictly
positive ACT/365F times \(t_i\), with at least one positive receipt. Their
present value is
\[
 P(r)=\sum_i\frac{c_i}{(1+r)^{t_i}},\qquad
 P'(r)=-\sum_i\frac{t_i c_i}{(1+r)^{t_i+1}}<0.
\]
Thus present value is continuous and strictly decreasing on \(r>-1\).
Certified endpoint values \(P(r_-)\geq100\) and \(P(r_+)\leq100\)
enclose the unique original rate. Mixed-sign receipts do not have this
guarantee and are outside this particular method. The bracket expresses
computational uncertainty about one rate, not different economic histories
or probabilities assigned to them.

The same rate must govern every later date. For elapsed days \(d\geq0\),
define the accumulation factor \(F(r,d)=(1+r)^{d/365}\). It increases with
\(r\), so
\[
 F(r_-,d)\leq F(r,d)\leq F(r_+,d).
\]
Each endpoint factor is itself enclosed by the integer-power calculation
used in the preceding dated method. Adding, subtracting and multiplying
these intervals, and taking reciprocals only of positive factors, encloses
the corresponding financial calculation. Starting with exact gross carrying
100 and allowance zero, the argument applies successively to accrual,
receipts, credit revisions, non-credit catch-ups and equal gross/allowance
write-offs. The original rate is never refitted when expected credit changes.

The POCI example makes the distinction concrete. Supply contractual receipts
of 60 and 75, separate from the initial expectations of 50 and 60. After
181 days, let \(f=x^{181/365}\). Reducing expected receipts to 40 and 48
reduces net carrying from \(100f\) to \(80f\), a loss of \(20f\).
A supported finding that one fifth of the entire claim has no reasonable
recovery prospect writes off gross and allowance of \(20f\) together.
Remaining gross and net are both \(80f\); the contractual terminal receipt
remaining after the write-off is 60. That contractual ceiling matters even
though the accounting write-off does not itself establish legal forgiveness
\citep[IFRS~9.5.4.4,B5.4.9]{eu20231803}.

At the first annual date, collect 40. The initial quadratic implies
\(80x-40=48/x\), so remaining gross is \(48/x\). If the supported expected
terminal receipt then improves to 60, net carrying is \(60/x\) and signed
allowance is \(-12/x\). Favorable changes in a POCI asset's lifetime loss
expectations can produce this negative allowance
\citep[IFRS~9.5.5.13--14]{eu20231803}. One year later, gross is 48, allowance
is \(-12\), and net is 60. Collecting 60 and clearing the remaining gross
and allowance closes the asset without another gain. Opening cash 300 becomes
\(300-98-2+40+60=300\); cumulative profit is zero.

For this calculation, a certified rate bracket of width \(2^{-40}\) gives
monetary intervals narrower than \(5.38\times10^{-10}\) currency units.
The profit interval contains zero, while cash is exactly 300. These are
bounds under the declared numerical budget. A wider rate interval can make
the financial result too imprecise; the method then reports that failure
instead of silently rounding a balanced result.

The returned bounds share the same original rate and cash history. Their
separate endpoints cannot be freely combined into a journal or a set of
jointly attainable statements. The documentary report therefore retains the
root certificate, its original cash-flow target, all dated origins and the
width of each amount. It proposes no postings from interval endpoints.
Direct remaining-cash present values and compounded gross origins check the
same history through different expressions. This extends the numerical
admission of the finite performing and POCI profiles; their supplied
classification, recovery, day-count and information-date requirements remain.
Floating rates, mixed-sign root selection and general modification or
write-off allocation methods require their own treatment.
